\documentclass[12pt, letterpaper]{article}
\usepackage[utf8]{inputenc}
\usepackage[T1]{fontenc}
\usepackage{times}
\usepackage[margin=1in]{geometry}
\usepackage{setspace}
\usepackage{parskip}
\usepackage{hyperref}
\usepackage{enumitem}
\usepackage{fancyhdr}
\usepackage{booktabs}
\usepackage{tabularx}
\usepackage{xcolor}
\usepackage{listings}
\usepackage{titlesec}
\usepackage{geometry}

\doublespacing
\setlength{\parindent}{0.5in}

\lstset{
  basicstyle=\footnotesize\ttfamily,
  breaklines=true,
  frame=single,
  numbers=left,
  numberstyle=\tiny,
  captionpos=t,
  aboveskip=12pt,
  belowskip=6pt,
  language=Python,
  keywordstyle=\bfseries\color{blue},
  commentstyle=\color{green!50!black},
  stringstyle=\color{red},
}

\pagestyle{fancy}
\fancyhf{}
\fancyhead[L]{\small\textsc{Arquitectura Multiagente -- Informe Técnico}}
\fancyhead[R]{\thepage}
\renewcommand{\headrulewidth}{0pt}

\title{\textbf{Informe Técnico: \\ Sistema Multiagente para Teoría del Caso Aumentada}}
\author{Santiago Esteban Redondo Perdomo \\ Juan David Conde Duarte}
\date{}

\begin{document}

\thispagestyle{empty}
\begin{center}
  \vspace*{2cm}
  {\LARGE\bfseries Informe Técnico: \\[6pt]
  Sistema Multiagente para Teoría del Caso Aumentada}\\[2cm]
  {\large Santiago Esteban Redondo Perdomo \\[4pt]
  Juan David Conde Duarte}\\[1cm]
  {\large \today}
  \vfill
\end{center}
\newpage

\section{Arquitectura General del Sistema}

El sistema sigue una arquitectura cliente-servidor. El backend está
implementado en Python 3 con Flask y expone tres endpoints REST:

\begin{itemize}[nosep]
  \item \texttt{POST /upload-pdf} -- Recibe un PDF, lo procesa
    mediante agentes y devuelve un JSON con entidades, matriz HPN,
    métricas y escenarios.
  \item \texttt{GET /network} -- Retorna el grafo multicapa
    previamente generado desde \texttt{outputs/network.json}.
  \item \texttt{GET /metrics} -- Retorna las métricas desde
    \texttt{outputs/metrics.json}.
  \item \texttt{GET /scenarios} -- Retorna los escenarios desde
    \texttt{outputs/scenarios.json}.
\end{itemize}

El frontend es una aplicación React 19 con Vite 8 que consume estos
endpoints y presenta un dashboard con ocho pestañas.

\subsection{Flujo de procesamiento de un PDF}

\noindent Al cargar un PDF, el backend ejecuta la siguiente pipeline:

\begin{enumerate}[nosep]
  \item \texttt{PDFService.extract(path)} $\rightarrow$
    \texttt{IngestionAgent.process\_pdf()} $\rightarrow$ \texttt{chunks[]}

    El \texttt{IngestionAgent} usa PyMuPDF (\texttt{fitz}) para
    extraer el texto página por página. Cada página se convierte en
    un chunk con la estructura:
    \texttt{\{page: int, text: str, hash: str\}}. El hash SHA-256
    permite detectar duplicados.

  \item \texttt{ExtractionService.execute(chunks)} $\rightarrow$
    \texttt{entities}

    Orquesta tres agentes de extracción:
    \texttt{FactualAgent}, \texttt{ActorsAgent} y
    \texttt{EvidenceAgent}.

  \item \texttt{HPNService.generate(chunks, entities)} $\rightarrow$
    \texttt{hpn\_matrix[]}

    Ejecuta el \texttt{NormativeAgent} para extraer normas del texto
    y luego el \texttt{HPNAgent} para cruzar hechos, pruebas y normas
    en una matriz estructurada.

  \item \texttt{ExportService.export\_hpn\_csv()},
    \texttt{export\_hpn\_json()}

    Persiste la matriz HPN en disco.

  \item \texttt{NetworkService.generate(hpn\_matrix, entities)}
    $\rightarrow$ \texttt{graph}

    El \texttt{NetworkAgent} construye un \texttt{networkx.DiGraph}
    con cinco capas de nodos y cuatro tipos de aristas.

  \item \texttt{ExportService.export\_graphml()},
    \texttt{export\_graph\_json()}

    Persiste el grafo en disco (GraphML y JSON).

  \item \texttt{MetricsService.generate(graph)} $\rightarrow$
    \texttt{metrics}

    El \texttt{MetricsAgent} calcula nodos, aristas, densidad,
    grado, betweenness centrality y PageRank.

  \item \texttt{ExportService.export\_metrics\_json()},
    \texttt{export\_metrics\_csv()}

    Persiste las métricas.

  \item \texttt{SimulationService.generate(graph)} $\rightarrow$
    \texttt{scenarios[]}

    El \texttt{SimulatorAgent} crea cuatro copias del grafo y
    aplica una mutación a cada una.

  \item \texttt{ExportService.export\_scenarios()}

    Persiste los escenarios.

  \item Se construye y retorna un JSON con:
    \texttt{\{success, pages, entities, hpn\_matrix, metrics,
    scenarios, summary\}}.
\end{enumerate}

\newpage
\section{Agentes de Extracción}

\subsection{FactualAgent}

\subsubsection*{Propósito}
Extraer hechos jurídicos del texto del PDF dividiendo el contenido
en oraciones y filtrando aquellas con longitud significativa.

\subsubsection*{Implementación}
\begin{itemize}[nosep]
  \item Usa el patrón \texttt{\textbackslash d\{1,2\}/\textbackslash
    d\{1,2\}/\textbackslash d\{4\}} para detectar fechas en el texto.
  \item Divide el texto de cada chunk por el carácter \texttt{.}
    para obtener oraciones candidatas.
  \item Filtra oraciones con menos de 30 caracteres (ruido).
  \item Cada hecho extraído contiene:
    \texttt{\{id, descripcion, fecha, fuente\_pagina\}}.
\end{itemize}

\subsubsection*{Limitación}
La división por puntos fragmenta incorrectamente oraciones con
abreviaturas comunes en documentos legales: ``S.A.S.'', ``No.'',
``art. 1602 C.C.''.

\subsection{ActorsAgent}

\subsubsection*{Propósito}
Identificar nombres de personas y partes procesales mencionadas en
el documento.

\subsubsection*{Implementación}
\begin{itemize}[nosep]
  \item Expresión regular:
    \texttt{[A-ZÁÉÍÓÚÑ][a-záéíóúñ]+(?:\textbackslash
    s[A-ZÁÉÍÓÚÑ][a-záéíóúñ]+)\{1,2\}}

    Captura secuencias de 2 a 3 palabras con inicial mayúscula
    (ej: ``Juan Pérez'', ``María García'').

  \item Lista de exclusión (\texttt{EXCLUDE\_WORDS}) con 35 términos
    legales: código, constitución, decreto, juzgado, tribunal,
    artículo, demanda, etc.

    Si alguna palabra del nombre candidato pertenece a esta lista,
    se descarta.

  \item Control de duplicados mediante un conjunto \texttt{found}.
\end{itemize}

\subsubsection*{Salida}
\texttt{\{nombre: str, tipo: str\}} con \texttt{tipo} fijo en
\texttt{"persona"}.

\subsection{EvidenceAgent}

\subsubsection*{Propósito}
Detectar en el texto las pruebas o medios de prueba mencionados.

\subsubsection*{Implementación}
\begin{itemize}[nosep]
  \item Lista fija de 9 keywords probatorias: contrato, correo,
    mensaje, factura, peritaje, audio, video, testimonio, documento.
  \item Búsqueda \texttt{in} (casefold) en el texto de cada chunk.
  \item Cada evidencia contiene:
    \texttt{\{id, descripcion, tipo, fuente\_pagina\}}.
\end{itemize}

\subsubsection*{Limitación}
Solo detecta la primera ocurrencia de cada keyword por chunk.
No distingue entre ``presentar un contrato'' y ``el contrato social
de la empresa''.

\subsection{NormativeAgent}

\subsubsection*{Propósito}
Extraer referencias normativas (leyes, artículos, códigos, decretos,
sentencias) del texto del documento.

\subsubsection*{Implementación}
\begin{itemize}[nosep]
  \item Lista de 8 keywords: ley, artículo, codigo, constitución,
    decreto, sentencia, jurisprudencia, norma.
  \item Para cada keyword, compila una expresión regular:
    \texttt{re.escape(keyword) + r"[\^{}.]\{0,120\}\."}

    Captura desde la keyword hasta el siguiente punto (máximo 120
    caracteres), obteniendo la oración completa donde aparece la
    referencia normativa.

  \item Control de duplicados con un conjunto \texttt{seen}
    (comparación casefold).
\end{itemize}

\subsubsection*{Salida}
\texttt{\{id, descripcion, pagina\}} donde
\texttt{descripcion} es la oración completa.

\newpage
\section{Matriz Hecho--Prueba--Norma (HPN)}

\subsection{Construcción}

La matriz HPN es la estructura de datos central del sistema. Se
construye en \texttt{HPNAgent.build\_hpn(facts, evidences, norms)}
mediante un cruce por índice secuencial:

\noindent Para cada hecho en la posición \texttt{i}:

\begin{itemize}[nosep]
  \item Se toma la prueba \texttt{evidences[i]} si existe;
    en caso contrario se asigna \texttt{"Sin prueba"}.
  \item Se toma la norma \texttt{norms[i]} si existe;
    en caso contrario se asigna \texttt{"Sin norma"}.
  \item El \texttt{estado\_epistemico} se define como:
    \begin{itemize}[nosep]
      \item \texttt{"Completo"} si existe tanto prueba como norma.
      \item \texttt{"Parcial"} si falta alguno de los dos.
    \end{itemize}
  \item El \texttt{riesgo} se define como:
    \begin{itemize}[nosep]
      \item \texttt{"Bajo"} si el estado es completo.
      \item \texttt{"Alto"} si el estado es parcial.
    \end{itemize}
\end{itemize}

\subsection{Estructura de cada fila}

\begin{lstlisting}[caption={Estructura de una fila de la matriz HPN}]
{
  "id":                "HPN-001",
  "elemento_juridico": "Pendiente Clasificacion",
  "hecho":             "El dia 15 de enero de 2025...",
  "prueba":            "contrato",
  "norma":             "articulo 1602 del Codigo Civil...",
  "fuente_expediente": "Pagina 1",
  "estado_epistemico": "Completo",
  "riesgo":            "Bajo",
  "contradicciones":   "Ninguna",
  "accion_sugerida":   "Validar soporte",
  "agente_responsable":"Agente HPN",
  "revision_humana":   "Pendiente"
}
\end{lstlisting}

\subsection{Utilidad jurídica}

La matriz HPN permite responder preguntas como:
\begin{itemize}[nosep]
  \item ¿Qué hechos tienen respaldo probatorio y normativo?
    (filas con estado \texttt{Completo} y riesgo \texttt{Bajo})
  \item ¿Qué hechos son débiles?
    (filas con estado \texttt{Parcial} o riesgo \texttt{Alto})
  \item ¿Qué normas están citadas y qué hechos activan?
  \item ¿Qué prueba respalda cada hecho?
\end{itemize}

\newpage
\section{Red Compleja Multicapa}

\subsection{Definición}

La red es un grafo dirigido (\texttt{networkx.DiGraph}) donde los
nodos representan entidades jurídicas organizadas en cinco capas y
las aristas representan relaciones semánticas entre ellas.

\subsection{Capas de nodos}

\begin{enumerate}[nosep]
  \item \textbf{Actores}: Nombres de personas extraídas por
    \texttt{ActorsAgent}. Atributo: \texttt{layer="actores"}.
  \item \textbf{Hechos}: Nodos con prefijo \texttt{FACT\_HPN-XXX}.
    Atributos: \texttt{layer="hechos"}, \texttt{label=hecho}.
  \item \textbf{Pruebas}: Nodos con prefijo \texttt{EVIDENCE\_HPN-XXX}.
    Atributos: \texttt{layer="pruebas"}, \texttt{label=prueba}.
  \item \textbf{Normas}: Nodos con prefijo \texttt{NORM\_HPN-XXX}.
    Atributos: \texttt{layer="normas"}, \texttt{label=norma}.
  \item \textbf{Riesgos}: Nodos con prefijo \texttt{RISK\_HPN-XXX}.
    Atributos: \texttt{layer="riesgos"}, \texttt{label=riesgo}.
\end{enumerate}

\subsection{Tipos de aristas}

\begin{itemize}[nosep]
  \item \textbf{soporta} ($EVIDENCE\_X \rightarrow FACT\_X$):
    La prueba X respalda al hecho X. Peso: 1.
  \item \textbf{activa} ($FACT\_X \rightarrow NORM\_X$):
    El hecho X activa la norma X. Peso: 1.
  \item \textbf{riesgo} ($RISK\_X \rightarrow FACT\_X$):
    El riesgo X está asociado al hecho X. Peso: 1.
  \item \textbf{participa} ($ACTOR \rightarrow FACT\_X$):
    El actor participa en el hecho X. Solo se crea si el nombre
    del actor aparece textualmente en la descripción del hecho
    (comparación \texttt{lower()}).
\end{itemize}

\subsection{Visualización en el frontend}

El componente \texttt{EnhancedNetworkGraph.jsx} renderiza el grafo
con \texttt{react-force-graph-2d} asignando colores por capa:

\begin{center}
\begin{tabular}{ll}
\toprule
\textbf{Capa} & \textbf{Color} \\
\midrule
Actores & \texttt{\#3b82f6} (azul) \\
Hechos & \texttt{\#f59e0b} (amarillo) \\
Pruebas & \texttt{\#10b981} (verde) \\
Normas & \texttt{\#8b5cf6} (púrpura) \\
Riesgos & \texttt{\#ef4444} (rojo) \\
\bottomrule
\end{tabular}
\end{center}

\noindent Las aristas también se colorean por tipo de relación:
soporta (verde), activa (púrpura), riesgo (rojo), participa (gris).
El nodo más central de cada capa se renderiza con mayor tamaño
(\texttt{nodeRelSize}). Incluye flechas direccionales y tooltip con
el ID del nodo.

\newpage
\section{Métricas Estratégicas}

\subsection{Métricas clásicas de red (backend)}

El \texttt{MetricsAgent} calcula las siguientes métricas sobre el
grafo usando \texttt{networkx}:

\begin{itemize}[nosep]
  \item \texttt{cantidad\_nodos}: Número total de nodos en el grafo.
  \item \texttt{cantidad\_aristas}: Número total de aristas.
  \item \texttt{densidad}: \texttt{nx.density(graph)}.
    $$D = \frac{2|E|}{|V|(|V|-1)}$$
  \item \texttt{grado}: \texttt{dict(graph.degree())}. Grado de cada
    nodo (suma de aristas entrantes y salientes).
  \item \texttt{betweenness}: \texttt{nx.betweenness\_centrality()}.
    Mide qué tan central es cada nodo como puente entre otros pares.
  \item \texttt{pagerank}: \texttt{nx.pagerank()}. Versión de
    PageRank adaptada a grafos dirigidos.
\end{itemize}

\subsection{Métricas estratégicas compuestas (frontend)}

El \texttt{StrategicMetricsPanel.jsx} calcula seis indicadores a
partir de la matriz HPN y las métricas de red:

\begin{enumerate}[nosep]
  \item \textbf{Cobertura}: $\frac{\text{filas completas}}{\text{total filas}} \times 100$

    Mide qué porcentaje de la matriz HPN tiene tanto prueba como
    norma asignadas.

  \item \textbf{Fragilidad}: $\frac{\text{filas con riesgo alto}}{\text{total filas}} \times 100$

    Porcentaje de elementos jurídicos en estado de alerta.

  \item \textbf{Centralidad}: $\frac{E}{N(N-1)} \times 100$

    Densidad relativa de la red. Indica qué tan interconectados
    están los elementos del caso.

  \item \textbf{Contradicción}: $\frac{\text{filas con contradicciones}}{\text{total filas}} \times 100$

    Porcentaje de filas donde se ha detectado alguna contradicción.

  \item \textbf{Robustez}: $100 - \text{fragilidad}$

    Porcentaje de elementos que no presentan riesgo alto.

  \item \textbf{Trazabilidad}: $\frac{\text{filas con prueba asignada}}{\text{total filas}} \times 100$

    Porcentaje de hechos que cuentan con respaldo probatorio.
\end{enumerate}

\noindent Cada métrica se muestra con un indicador de semáforo:
verde (bueno), amarillo (regular), rojo (crítico), según umbrales
definidos en el frontend.

\newpage
\section{Simulación de Escenarios}

\subsection{Implementación}

El \texttt{SimulatorAgent.simulate(graph)} recibe el grafo base
(estado actual del caso) y genera cuatro escenarios hipotéticos
mediante mutaciones controladas del grafo:

\begin{enumerate}[nosep]
  \item \textbf{Exclusión de prueba}: Clona el grafo y elimina la
    primera arista de la lista (\texttt{list(graph.edges())[0]}).
    Reduce el número de aristas en 1.

    \textit{Utilidad}: Modela qué sucede si se declara inadmisible
    o se pierde una prueba durante el proceso.

  \item \textbf{Nuevo peritaje}: Clona el grafo y agrega un nodo
    \texttt{NEW\_EXPER} en la capa \texttt{pruebas}. Incrementa el
    número de nodos en 1.

    \textit{Utilidad}: Evalúa el impacto de incorporar un dictamen
    pericial adicional.

  \item \textbf{Testigo contradictorio}: Clona el grafo y agrega un
    nodo \texttt{CONTRADICTORY\_WITNESS} en la capa \texttt{pruebas}.
    Incrementa el número de nodos en 1.

    \textit{Utilidad}: Modela la aparición de un testigo que
    contradice la versión de los hechos.

  \item \textbf{Nueva norma favorable}: Clona el grafo y agrega un
    nodo \texttt{NEW\_NORM} en la capa \texttt{normas}. Incrementa
    el número de nodos en 1.

    \textit{Utilidad}: Simula el efecto de invocar una norma
    adicional o una jurisprudencia favorable.
\end{enumerate}

\subsection{Salida}

\begin{lstlisting}[caption={Ejemplo de respuesta del simulador}]
[
  {"name": "Exclusion de prueba",   "nodes": 101, "edges": 95},
  {"name": "Nuevo peritaje",        "nodes": 102, "edges": 96},
  {"name": "Testigo contradictorio","nodes": 102, "edges": 96},
  {"name": "Nueva norma favorable", "nodes": 102, "edges": 96}
]
\end{lstlisting}

\subsection{Visualización en el dashboard}

El \texttt{AdvancedScenariosPanel.jsx} muestra los escenarios en un
grid de tarjetas. La primera tarjeta corresponde al estado actual
(base) y las siguientes a cada escenario simulado. Para cada
escenario se calcula el delta ($\Delta$) de nodos y aristas respecto
a la base, con color verde para deltas positivos y rojo para
negativos. Se incluye una etiqueta textual del impacto (positivo,
negativo, neutro).

\newpage
\section{Componentes del Dashboard}

El frontend organiza la información en ocho pestañas, cada una
implementada como un componente React independiente.

\subsection{Vista General}
\label{sec:vista-general}

\texttt{Componente:} definido inline en \texttt{Dashboard.jsx}
(ReadinessBar, ActionableInsights, CaseOverview).

\begin{itemize}[nosep]
  \item \textbf{Barra de preparación del caso} (\texttt{ReadinessBar}):
    Barra horizontal segmentada que muestra la proporción de filas
    completas (verde), parciales (amarillo) y críticas (rojo) en la
    matriz HPN. Incluye etiquetas con conteos absolutos.

  \item \textbf{Señales accionables} (\texttt{ActionableInsights}):
    Siete tarjetas con análisis interpretativo generado
    automáticamente:
    \begin{itemize}[nosep]
      \item Hechos afirmables con fuerza (completos + riesgo bajo)
      \item Hechos a modular (parciales o riesgo alto)
      \item Prueba crítica a proteger (tipo de prueba único)
      \item Documentos a preparar primero (orden de tipos probatorios)
      \item Excepción a anticipar (filas con contradicciones)
      \item Puntos no aptos para memorial (baja trazabilidad)
    \end{itemize}

  \item \textbf{Entidades del caso} (\texttt{CaseOverview}):
    Grid de cuatro tarjetas mostrando: Actores, Hechos Jurídicos,
    Pruebas y Normas Citadas. Cada tarjeta lista las entidades
    extraídas con un badge identificador.
\end{itemize}

\subsection{Matriz HPN}

\texttt{Componente: FilterableHPNTable.jsx}

\begin{itemize}[nosep]
  \item Tabla HTML con las columnas: ID, Hecho, Prueba, Norma,
    Estado, Riesgo, Contradicciones, Acción Sugerida.
  \item \textbf{Filtros interactivos} (seis controles):
    \begin{itemize}[nosep]
      \item Estado epistémico (desplegable con valores únicos)
      \item Riesgo (desplegable)
      \item Prueba (desplegable)
      \item Norma (desplegable)
      \item Acción sugerida (desplegable)
      \item Búsqueda textual (campo de texto, busca en todas las
        columnas)
    \end{itemize}
  \item Botón ``Limpiar'' para resetear todos los filtros.
  \item Conteo de resultados: ``Mostrando X de Y elementos''.
  \item Filas coloreadas: rojo para riesgo alto, amarillo para
    estado parcial.
  \item Badges de estado y riesgo con estilos diferenciados.
\end{itemize}

\subsection{Red Multicapa}

\texttt{Componente: EnhancedNetworkGraph.jsx}

\begin{itemize}[nosep]
  \item Grafo interactivo usando \texttt{react-force-graph-2d}.
  \item \textbf{Leyenda visual}: Muestra las cinco capas con sus
    colores correspondientes y los cuatro tipos de aristas.
  \item Nodos coloreados por capa, con tamaño variable según tipo
    (actores más grandes).
  \item Aristas coloreadas por tipo de relación, con flechas
    direccionales.
  \item Tooltip al pasar el mouse: muestra el ID del nodo.
  \item \texttt{nodeLabel}: muestra el ID como etiqueta.
  \item Dimensiones: 1200x600px, con scroll si es necesario.
\end{itemize}

\subsection{Métricas Estratégicas}

\texttt{Componente: StrategicMetricsPanel.jsx}

\begin{itemize}[nosep]
  \item Grid de seis tarjetas, cada una con:
    \begin{itemize}[nosep]
      \item Nombre de la métrica
      \item Valor numérico con porcentaje
      \item Descripción corta
      \item Borde izquierdo coloreado (semáforo verde/amarillo/rojo)
    \end{itemize}
  \item Las seis métricas: Cobertura, Fragilidad, Centralidad,
    Contradicción, Robustez, Trazabilidad.
  \item Subtítulo: ``Priorización basada en la matriz HPN y la
    estructura de red multicapa.''
\end{itemize}

\subsection{Alertas}

\texttt{Componente: EnhancedAlertsPanel.jsx}

\begin{itemize}[nosep]
  \item Cinco grupos de alertas, cada uno con un color distintivo:
    \begin{itemize}[nosep]
      \item \textbf{Vacíos críticos -- Alto riesgo}: Filas con
        riesgo \texttt{"Alto"}. Fondo rojo.
      \item \textbf{Pruebas únicas -- Críticas de proteger}: Tipos
        de prueba que aparecen una sola vez en la matriz. Fondo azul.
      \item \textbf{Contradicciones detectadas}: Filas donde
        \texttt{contradicciones} es distinto de \texttt{"Ninguna"}.
        Fondo rosa.
      \item \textbf{Baja trazabilidad}: Filas con estado parcial
        y sin prueba asignada. Fondo púrpura.
      \item \textbf{Sin soporte probatorio}: Filas sin prueba
        asociada (\texttt{"Sin prueba"} o vacío). Fondo azul.
    \end{itemize}
  \item Cada alerta muestra el ID de la fila y un extracto del hecho.
  \item Limitación a 10 elementos por grupo para evitar saturación
    visual.
\end{itemize}

\subsection{Simulador}

\texttt{Componente: AdvancedScenariosPanel.jsx}

\begin{itemize}[nosep]
  \item Grid de tarjetas. La primera es el estado actual (base),
    resaltada con borde azul.
  \item Cada escenario muestra:
    \begin{itemize}[nosep]
      \item Nombre del escenario
      \item Nodos y aristas absolutos
      \item $\Delta$ nodos y $\Delta$ aristas respecto a la base
        (con color según signo)
      \item Etiqueta de impacto: positivo (verde), negativo (rojo),
        neutro (gris)
    \end{itemize}
  \item Etiquetas de impacto predefinidas para cada escenario:
    \begin{itemize}[nosep]
      \item Exclusión de prueba: Negativo
      \item Nuevo peritaje: Positivo
      \item Testigo contradictorio: Positivo
      \item Nueva norma favorable: Positivo
    \end{itemize}
\end{itemize}

\subsection{Preguntas Sugeridas}

\texttt{Componente: QuestionsPanel.jsx}

\begin{itemize}[nosep]
  \item Genera preguntas automáticas basadas en la matriz HPN y las
    entidades extraídas.
  \item Cuatro categorías, cada una con un color de borde:
    \begin{itemize}[nosep]
      \item \textbf{Interrogatorio directo} (borde púrpura):
        Preguntas sobre hechos completos y sin controversia.
      \item \textbf{Contraexamen} (borde amarillo): Preguntas sobre
        hechos débiles, controvertidos o con contradicciones.
      \item \textbf{Aclaración al cliente} (borde verde): Preguntas
        sobre tipos de prueba y disponibilidad de documentos.
      \item \textbf{Consulta a perito} (borde rosa): Preguntas sobre
        debilidades probatorias que requieren peritaje.
    \end{itemize}
  \item Las preguntas se construyen interpolando datos reales del
    caso: IDs de fila HPN, fragmentos de hechos, tipos de prueba,
    nombres de actores.
  \item Agrupación por categoría con subtítulos y conteos.
\end{itemize}

\subsection{Reporte Exportable}

\texttt{Componente: ExportPanel.jsx}

\begin{itemize}[nosep]
  \item Genera un documento HTML completo con:
    \begin{itemize}[nosep]
      \item Resumen ejecutivo (metadatos del caso, conteo de
        entidades)
      \item Matriz HPN completa en tabla HTML
      \item Entidades del caso (actores, hechos, pruebas, normas)
      \item Métricas estratégicas (cobertura, fragilidad,
        trazabilidad)
      \item Escenarios simulados
    \end{itemize}
  \item \textbf{Botón ``Exportar PDF''}: Abre el HTML en una nueva
    ventana y ejecuta \texttt{window.print()}, permitiendo al
    usuario guardar como PDF desde el diálogo de impresión del
    navegador.
  \item \textbf{Botón ``Exportar HTML''}: Descarga el documento como
    archivo \texttt{.html} mediante \texttt{Blob} y
    \texttt{URL.createObjectURL()}.
  \item Vista previa del reporte dentro del dashboard, renderizada
    con \texttt{dangerouslySetInnerHTML}.
  \item Estilos CSS embebidos en el HTML para garantizar formato
    consistente al imprimir.
\end{itemize}

\newpage
\section{Diagrama de Arquitectura}

\noindent El siguiente diagrama de flujo muestra el recorrido
completo de los datos desde la carga del PDF hasta la visualización
en el dashboard:

\begin{verbatim}
+-----------+     +------------------+     +------------------+
|   PDF     | --> | IngestionAgent   | --> | ExtractionService|
| (upload)  |     | (fitz) chunks[]  |     |                   |
+-----------+     +------------------+     +--+----+----+------+
                                               |    |    |
                          FactualAgent <--------+    |    |
                          ActorsAgent  <-------------+    |
                          EvidenceAgent<------------------+
                                               |
                                               v
                                     +------------------+
                                     | NormativeAgent   |
                                     | (regex oraciones)|
                                     +--------+---------+
                                              |
                                              v
                                  +-----------------------+
                                  | HPNAgent              |
                                  | (cruza por índice)    |
                                  +----------+------------+
                                             |
                                   +---------+---------+
                                   |                   |
                                   v                   v
                            +-----------+     +------------------+
                            | Export    |     | NetworkAgent     |
                            | CSV/JSON  |     | (networkx.DiGraph)|
                            +-----------+     +--------+---------+
                                                      |
                                            +---------+---------+
                                            |                   |
                                            v                   v
                                     +-----------+     +------------------+
                                     | Metrics   |     | SimulatorAgent   |
                                     | (nx.grado,|     | (4 escenarios)   |
                                     |  between, |     +--------+---------+
                                     |  pagerank)|              |
                                     +-----+-----+              |
                                           |                    |
                                           v                    v
                                     +-----------+     +------------------+
                                     | Export    |     | Export           |
                                     | JSON/CSV  |     | scenarios.json   |
                                     +-----------+     +------------------+
                                                               
                                                               
  +--------------------------------------------------------------+
  |                    RESPONSE JSON                              |
  | {pages, entities, hpn_matrix, metrics, scenarios, summary}   |
  +--------------------------------------------------------------+
                           |
                           v
  +--------------------------------------------------------------+
  |                    REACT DASHBOARD                            |
  |  +-------------+  +-------------+  +--------------------+    |
  |  | Vista       |  | Matriz HPN  |  | Red Multicapa      |    |
  |  | General     |  | (filtros)   |  | (force-graph)      |    |
  |  +-------------+  +-------------+  +--------------------+    |
  |  +-------------+  +-------------+  +--------------------+    |
  |  | Métricas    |  | Alertas     |  | Simulador          |    |
  |  +-------------+  +-------------+  +--------------------+    |
  |  +-------------+  +-------------+                             |
  |  | Preguntas   |  | Reporte     |                             |
  |  | Sugeridas   |  | Exportable  |                             |
  |  +-------------+  +-------------+                             |
  +--------------------------------------------------------------+
\end{verbatim}

\section{Tecnologías Utilizadas}

\begin{center}
\begin{tabularx}{\textwidth}{lX}
\toprule
\textbf{Componente} & \textbf{Tecnología} \\
\midrule
Backend framework & Flask 3.0 (Python 3.10+) \\
Extracción de texto PDF & PyMuPDF (fitz) \\
Procesamiento de grafos & NetworkX 3.0 \\
Exportación de datos & pandas, JSON, GraphML \\
Frontend framework & React 19 \\
Herramienta de build & Vite 8 \\
Grafo interactivo & react-force-graph-2d 1.29 \\
Peticiones HTTP & Axios 1.18 \\
Linter & oxlint 1.69 \\
\bottomrule
\end{tabularx}
\end{center}

\end{document}
